\subsection{十三个算法与模型学习单元}

\setlength{\tabcolsep}{2.5pt}
\small
\begin{longtable}{p{1.25cm}p{2.55cm}p{6.25cm}p{5.45cm}}
\caption{成员学习内容、代码逻辑与使用边界}\label{tab:learning-units}\\
\toprule
成员 & 算法或模型 & 基本原理、步骤与代码逻辑 & 适用问题、优缺点和限制\\
\midrule
\endfirsthead
\toprule
成员 & 算法或模型 & 基本原理、步骤与代码逻辑 & 适用问题、优缺点和限制\\
\midrule
\endhead
\bottomrule
\endfoot

队长 & ODE与RK4
& 先从受力、守恒或状态转移关系写出 $\dot{y}=f(t,y)$，，再由初值出发，用四个斜率的加权值向前推进。
代码入口为 \texttt{rhs(t,y,p)} 和 \texttt{rk4\_step(...)}；用Euler法和
\texttt{solve\_ivp}/\texttt{ode45}作对照，并保存步长误差。
& 适合状态随时间连续变化的力学、传热、种群和传播过程。RK4在精度与实现难度之间较为均衡，
但刚性方程会受稳定域限制；此时应改用刚性求解器，不靠缩小一步掩盖问题。\\

队长 & PDE与有限差分
& 空间区域、初始条件和边界条件先落到网格上，再将时间、空间导数换成差分格式。程序依次生成
网格、装配差分矩阵、推进时间层，并比较三档网格的同一观测量。
& 适合规则区域上的热、波和扩散问题。公式与网格位置对应直观，代价是复杂边界难处理，
显式格式还须核对稳定条件。\\

队长 & 有限元模型
& 这部分从PDE的弱形式入手：划分单元，计算局部矩阵，装配全局稀疏矩阵，施加
边界后求节点量。代码模块依次为 \texttt{mesh}、\texttt{element}、\texttt{assemble}、
\texttt{bc} 和 \texttt{solve}。
& 适合不规则几何、变系数介质和混合边界。几何适应性强，但网格质量、形函数和边界
装配都容易引入错误；规则区域的入门题先用有限差分作基线。\\

队长 & 几何递推与事件判定
& 已知龙头位置后，先列相邻把手的定长方程；筛去不满足构件次序和运动方向的候选根，随后逐节
递推。碰撞用矩形顶点和分离轴判定，时间推进遇到第一次触发即停。代码保留
\texttt{next\_handle}、\texttt{collision} 和状态码。
& 适合链式刚体、轨迹跟随和首达事件。优点是公式、程序循环和物理对象能对读；缺点是
多根、步长和漏检风险要另列检查，不能只写“取合理解”。\\

组员1 & LP与MILP
& 先逐项列出决策变量、线性目标、等式和不等式；遇到启停、选择或分配，才引入0--1变量。程序先算出
小规模枚举答案，再调用PuLP/Pyomo或MATLAB求解器，读取可行性、目标值和最优间隙。
& 适合资源分配、排程、选址和生产决策。线性模型便于解释并能给最优性信息；整数规模
大时计算量上升，Big--M须来自实际上界，不能任意指定过大的常数。\\

组员1 & 多目标与NSGA--II
& 各项目标先分别保留，不急着揉成一个总分。染色体经可行性修复后，再按非支配排序和拥挤
距离选择，输出Pareto解集。代码记录种群、代数、随机种子和每个候选的原始指标。
& 适合成本、收益、风险等目标互相冲突的问题。它能展示取舍，但不替参赛者选最终方案；
变量少且约束凸时，$\varepsilon$约束法可直接呈现取舍边界。\\

组员1 & 粒子群优化
& 每个粒子都保留位置、速度、个体最好解和群体最好解；迭代时包含速度更新、越界处理和重新
评价。代码把目标函数、约束罚项、种子和停止阈值分开，至少做多次独立运行。
& 适合连续、非光滑或难求导的参数搜索，编码简短；它没有连续全局最优证书，结果会受
边界、罚项和随机种子影响。若变量只有一维且单调，可改用遍历或二分。\\

组员1 & 图论与网络优化
& 建模时把对象落成节点、关系落成带权边。非负单源最短路用Dijkstra；容量与费用同时出现时
用最小费用流。代码先建图，再运行算法，最后把路径或流量映射回现实对象并核对容量。
& 适合运输、路径、匹配和网络资源配置。结构清楚、结果易复算；动态图、负权边和多目标
路线要另行处理，最小生成树也不能代替两点间最短路。\\

组员2 & 缺失、异常与重复处理
& 拿到附件后先建数据字典和主键，并把未检出、未知、未采样分开记录；随后才按成因决定删除、插值或KNN，
异常值用箱线图、$3\sigma$和业务范围复核。代码以训练集拟合清洗参数并输出处理记录。
& 适合一切附件数据题。它决定后续样本是否可信，但不能把业务极值自动删除，也不能在
划分训练集前用全体数据估计均值或尺度。\\

组员2 & PCA与K--Means
& 数据标准化后，对协方差矩阵作特征分解并按解释率选取主成分；降维结果再送入K--Means，
比较不同簇数的轮廓系数。代码输出载荷、解释率、簇中心和原样本编号。
& PCA适合线性降维，K--Means适合近似球形簇。两者速度快、图形直观，但主成分未必有
直接业务含义，密度不均或非凸簇应比较DBSCAN或谱聚类。\\

组员2 & 熵权与TOPSIS
& 指标须先正向化和无量纲化。熵权刻画样本差异，TOPSIS据此计算对象到正、负理想解的距离。
代码逐步导出标准矩阵、权重、距离和排序，并对权重作扰动。
& 适合多指标评价与方案排序。计算过程清楚，但“差异大”不等于“业务上更重要”；指标
方向、相关重复和权重敏感性都要说明，不能只给最终名次。\\

组员2 & 回归与分类基线
& 连续响应以线性、Ridge和Lasso为基线；类别响应以Logistic为基线。保持同一训练/验证划分，再与
随机森林或XGBoost比较。代码用pipeline串起预处理、拟合、预测和指标，保存系数或特征贡献。
& 适合关系估计、风险评分和类别判断。线性模型可解释，树模型能处理非线性；后者容易
过拟合且外推较弱。相关或高准确率都不构成因果证据。\\

组员2 & ARIMA与滚动验证
& 先画出序列，辨认趋势和季节性，再检查差分后的平稳性并选择ARIMA/SARIMA阶数。每次预测只读取当前
时点以前的数据，按滚动窗口计算MAE、RMSE和MAPE。代码保存每个预测起点和误差。
& 适合有时间顺序的需求、价格和状态预测。统计模型对中小样本友好，但结构突变会使历史
规律失效；时间序列不能随机打乱做普通K折。\\
\end{longtable}
\normalsize

